\subsection{Source locality removes artificial cutoff amplification}
The global bound \eqref{eq:jetallowances} loses the source decay before
applying the heat weight. A cellwise replacement preserves it. Let
$[a_i,b_i]$ be a source cell of length $h_i$, let $B_i$ bound
$\sup_{[a_i,b_i]}|\Phi''|$, and let $\lambda_i$ be the nodal hats.
For implemented nodal data $v_i$, assume
$|v_i-\Phi_N(u_i)|\le\sigma_i$, with nonnegative interpolant $f$.
Define the spatially dependent omitted-mode allowance
\begin{equation}
 \tau_N(v)=\frac{2\pi^2(N+1)^4e^{9v-\pi(N+1)^2e^{4v}}}
 {1-(1+1/(N+1))^4e^{-\pi(2N+3)e^{4v}}}.
 \label{eq:localmodetail}
\end{equation}
For a time ceiling $b$, write $W_j(u)=u^je^{bu^2+Yu}$ and retain
$T_j$ from \eqref{eq:jetallowances} with $s=\max\{b,0\}$.

\begin{proposition}[Cellwise jet comparison]\label{prop:localjets}
Uniformly for $t\le b$ in the chosen slab, all real $x$, and $|y|\le Y$,
\begin{align}
 |\partial_z^j(H-H_f)|&\le\delta_j^{\rm loc},\qquad 0\le j\le3,\notag\\
 \delta_j^{\rm loc}
 &=\frac12\sum_i B_i\int_{a_i}^{b_i}(u-a_i)(b_i-u)W_j(u)\,du\notag\\
 &\quad+\sum_i\bigl(\tau_N(u_i)+\sigma_i\bigr)
                \int_0^U\lambda_i(u)W_j(u)\,du+T_j.
 \label{eq:localjets}
\end{align}
For the complete-source cell constants below, the interpolation term is
no larger than $B_2^{\Phi}h_s^2F_j/8$. For fixed $b,Y,j$ and maximum cell length
bounded above, its coefficient of $h_s^2$ remains bounded as $U$ grows.
\end{proposition}
\begin{proof}
Linear interpolation of the complete source costs
$B_i(u-a_i)(b_i-u)/2$. The positive hats transfer the nodal mode and
arithmetic errors, while consecutive omitted-mode ratios prove
\eqref{eq:localmodetail}. Integrating these bounds against $W_j$ and
adding the complete source tail proves \eqref{eq:localjets}.

To specify $B_i$, put $A_n=\pi n^2e^{4a_i}$ and $B_n=\pi n^2e^{4b_i}$.
Assign coefficients $75,330,224,32$ to the powers $5/4,9/4,13/4,17/4$,
respectively. Write $m_p(A,B)=v^pe^{-v}$, where
$v=\min\{B,\max\{A,p\}\}$. Then
\[
 B_i=\pi^{-1/4}\sum_{n\ge1}n^{-1/2}\sum_p c_p m_p(A_n,B_n)
 \le B_2^{\Phi}.
\]
This follows by bounding the four monomials in the differentiated
source as in Lemma~\ref{lem:B2}. For a finite sum through $N_B\ge1$,
the omitted tail is at most its first omitted expression at $a_i$
divided by
$1-(1+1/(N_B+1))^8e^{-\pi(2N_B+3)e^{4a_i}}$.
Each omitted monomial is decreasing in $u$.

For cell lengths at most $h_0$, a bound for $B_i$ at $u$ is the
supremum of this positive derivative envelope over
$|v-u|\le h_0$, $v\ge0$. Its far tail is bounded by a polynomial in
$e^{4(u+h_0)}$ times $e^{-\pi e^{4(u-h_0)}}$. This is integrable
against $W_j$. Using $(u-a_i)(b_i-u)\le h_s^2/4$ proves the final
assertion. The original coefficient $B_2^{\Phi}F_j/8$, in contrast, diverges
with $U$ when $b>0$. This distinction concerns artificial cutoff
amplification, not the genuine time growth of the approximation cost.
\end{proof}
A quadrature-free upper bound for each interpolation contribution is
$B_i h_i^3\max_{[a_i,b_i]}W_j/12$; the nodal contribution is at most
$h_i\max W_j$ times the mean of its two endpoint allowances.
The positive maxima occur at endpoints or positive solutions of
$2bu^2+Yu+j=0$. Taking the minimum with the original bound preserves
dominance even on a coarse mesh. Nodal arithmetic retains the source
decay only when its pointwise allowances do; a relative rounding
hypothesis must be verified. Note~23 \cite{n23} records the separate
interpolation, mode, tail and arithmetic terms.

\subsection{Local amplitude improves nonlinear transfer}
At a comparison point put $g=H_f$, $g_j=\partial_z^jH_f$, and let
$d_j$ be any valid jet allowances. Set
\[
 P_j=|g_j|+d_j,\quad m=(|g|-d_0)_+,\quad M=|g|+d_0,\quad q=\epsilon_g^2.
\]
Define maxima over $m\le v\le M$:
\begin{align}
 A_*&=\frac{q}{\pi(q+m^2)^2},&
 A_1&=\max_v\frac{4qv}{\pi(q+v^2)^3},\notag\\
 a_*&=\max_v\frac{v}{q+v^2},& b_*&=\frac1{q+m^2},&
 w_1&=\max_v\frac{2qv}{(q+v^2)^2}.\label{eq:localrational}
\end{align}
The three maxima are evaluated at the clamped points
$\epsilon_g/\sqrt5$, $\epsilon_g$, and $\epsilon_g/\sqrt3$.
Set
\begin{align}
 e_\rho&=2A_*P_1d_1+A_1P_1^2d_0,\notag\\
 e_j&=A_*(P_1d_2+P_2d_1)+A_1P_1P_2d_0,\notag\\
 e_{13}&=A_*(P_1d_3+P_3d_1)+A_1P_1P_3d_0,\notag\\
 e_w&=w_1d_0,\qquad e_S=2a_*d_0,\qquad \rho_*=A_*P_1^2,\notag\\
 e_{\rho,t}&=2e_{13}+|r_t|(e_\rho+2\rho_*e_w)
 +4\{P_2a_*e_\rho+\rho_*(a_*d_2+P_2b_*d_0)\}.
 \label{eq:localnonlinear}
\end{align}

\begin{proposition}[Amplitude-aware comparison]\label{prop:localnonlinear}
The quantities in \eqref{eq:localnonlinear} bound the corresponding
actual density, current, potential and density-time-derivative errors.
In addition,
\begin{align}
 |S_t-(S_f)_t|&\le2(a_*d_2+P_2b_*d_0)+|r_t|e_w,\notag\\
 |S_i|,|(S_f)_i|&\le G_i:=2P_1a_*+|r_i|,\notag\\
 |S_i-(S_f)_i|&\le d_i^S:=2a_*d_1+2P_1b_*d_0+|r_i|e_w.
 \label{eq:localpotential}
\end{align}
No simple-root or nonvanishing-field premise is needed.
\end{proposition}
\begin{proof}
The segment between the two field values lies in the annulus
$m\le|h|\le M$. The constants \eqref{eq:localrational} bound
$A(h)$, $|\nabla_h A|$, $|h/(|h|^2+q)|$, the derivative norm of
that map, and $|\nabla_h w|$, respectively.
For the fourth bound the radial eigenvalues are
$(q-v^2)/(q+v^2)^2$ and $1/(q+v^2)$.
The product comparison and exact time-derivative argument from
Proposition~\ref{prop:fieldstability} now give
\eqref{eq:localnonlinear}. The exact formulas for $S_t,S_i$ give
\eqref{eq:localpotential} directly.
All constants are bounded by their global counterparts in
\eqref{eq:fieldstability}--\eqref{eq:gradientstability}.
For example, relative to $1/(\pi\epsilon_g^2)$, $A_*$ gains the
factor $[1+(m/\epsilon_g)^2]^{-2}$.
The gain may disappear at a zero; the bound remains valid there.
\end{proof}
No physical normalization has changed. In particular there is no
unpaid time derivative of a rescaled field.

\subsection{Local covariance and evaluation budgets}
The local envelopes may be filtered before taking suprema.
For any pointwise error $e_V$,
$|\partial^\alpha K*(V-V_f)|\le|\partial^\alpha K|*e_V$.
With $d_i=d_i^S$ from \eqref{eq:localpotential}, direct comparison of
the actual covariances gives
\begin{align}
 |Q-Q_f|&\le2\sum_i\{K*(G_id_i)+(K*G_i)(K*d_i)\},\notag\\
 |\Delta(Q-Q_f)|&\le\sum_i\biggl\{2|\Delta K|*(G_id_i)
 +2(K*d_i)(|\Delta K|*G_i)\notag\\
 &\qquad+2(K*G_i)(|\Delta K|*d_i)
 +4\sum_{a=x,y}(|K_a|*G_i)(|K_a|*d_i)\biggr\}.
 \label{eq:localcovariance}
\end{align}
To prove this, move the derivatives of the averaged squares onto $K$
and use $\Delta(v^2)=2v\Delta v+2|\nabla v|^2$ for each averaged
gradient. This is the proof of Theorem~\ref{thm:covariance} with its
spatial information retained. For constant envelopes the Laplacian
coefficient becomes
$6\|\Delta K\|_1+4(\|K_x\|_1^2+\|K_y\|_1^2)$, slightly sharper
than \eqref{eq:covariancestability}. Positive convolution bounds still
need a numerical integration enclosure if evaluated by quadrature.

There is also a finite evaluation formula without division by $t$ or $z$.
On a cell centered at $c$ with half-width $d$, expand
\[
 e^{tu^2\pm izu}=e^{tc^2\pm izc}
 \sum_{n=0}^{M_e}\frac{[(2tc\pm iz)(u-c)+t(u-c)^2]^n}{n!}
       +\text{remainder}.
\]
Multiply by the linear source and $u^j$, and integrate the polynomial
exactly, with derivative factors $(\pm i)^j/2$ for the two branches.
Writing $b=c+d$ and $\Theta=(2|t|c+|z|)d+|t|d^2$, the cell's
absolute jet truncation error is at most
\begin{equation}
 2d\max\{f(c-d),f(c+d)\}b^j
 e^{tc^2+|y|c+\Theta}\frac{\Theta^{M_e+1}}{(M_e+1)!}.
 \label{eq:localevaluation}
\end{equation}
This follows from the absolute exponential remainder. Sum over cells.
It pays truncation, not rounding. The exact $H_f$ still solves heat;
a truncated implementation is compared to it with this additional error.
The raw error-function subtraction can be unusable near its removable
singularity even at high precision; Note~23 records a failed case and
the separate local evaluation check.

For the first trial, Note~23 fixes $[-16,64]$, observation rectangle
$[-8,8]\times[-1/2,1/2]$, the full buffer
$[-10,10]\times[-1,1]$, widths $(2,1/2)$ and $\eta=1/2$.
Its absolute contract includes $R,J,R_t,U,Q,\Delta Q$, the required
filtered derivatives and source-derived boundary traces.
On this slab $r_x=r_y=0$, so the spatial curvature correction is exactly
zero, while the time forcing and current source remain.
The 192-cell diagnostic does not meet the contract. The proved local
bounds are an improvement to the comparison method, not an informative
full-transition representation or a certified numerical solution.
